% ============================================================================
% Chuong 4: CAI DAT, TRIEN KHAI VA THU NGHIEM HE THONG
% ============================================================================

\chapter{CÀI ĐẶT, TRIỂN KHAI VÀ THỬ NGHIỆM HỆ THỐNG}
\label{ch:impl}

Chương này trình bày \emph{toàn bộ} những gì đã được hiện thực và vận hành: từ hạ
tầng nền tảng, qua ba thành phần lõi của vòng phát hiện~--~phản hồi (bộ thu thập
eBPF, bộ máy pipeline và trung tâm vận hành an ninh~--~SOC), tới ứng dụng giao diện
cho người vận hành, và cuối cùng là việc chạy thử nghiệm end-to-end trên cụm
Kubernetes thật. Theo đúng trọng tâm của đề tài, phần trình bày \emph{đào sâu} vào
bộ thu thập, pipeline và SOC~--~nơi đặt thuật toán phát hiện và logic phản hồi~--~kèm
sơ đồ luồng, mã giả và trích đoạn mã nguồn/cấu hình thật; phần hạ tầng và triển khai
chỉ được \emph{giới thiệu vắn tắt} nhưng vẫn liệt kê đủ thành phần. Những đoạn mã dài
hơn được đưa vào phần Phụ lục. Phần đánh giá
\emph{định lượng} (tỷ lệ phát hiện, FPR, chi phí, độ trễ) tách riêng sang
Chương~\ref{ch:eval}; còn thiết kế kiến trúc/thành phần (sơ đồ) đã trình bày ở
Chương~\ref{ch:design} và sẽ được tham chiếu lại khi cần.

\section{Tổng quan hiện thực}
\label{sec:impl-overview}

Hệ thống được tổ chức thành nhiều kho mã nguồn độc lập, mỗi kho phụ trách một tầng và
được đóng gói thành ảnh container riêng để xây dựng, kiểm thử, nâng cấp tách bạch.
Bảng~\ref{tab:impl-stack} ánh xạ mỗi thành phần tới kho mã, ngôn ngữ, đối tượng
Kubernetes và tệp mã nguồn cốt lõi; đây cũng là bản đồ tra cứu cho toàn bộ chương.

\begin{table}[H]
  \centering
  \caption{Bản đồ hiện thực: thành phần $\rightarrow$ kho mã, ngôn ngữ, đóng gói và tệp cốt lõi.}
  \label{tab:impl-stack}
  \scriptsize
  \begin{tabularx}{\linewidth}{@{}l l l l X@{}}
    \toprule
    \textbf{Thành phần} & \textbf{Kho mã} & \textbf{Ngôn ngữ} & \textbf{Đối tượng K8s} & \textbf{Tệp cốt lõi} \\
    \midrule
    Bộ thu thập eBPF      & \texttt{collector} & C (nhân) + C++ & \texttt{DaemonSet} (đặc quyền) & \texttt{tcp\_connect.bpf.c}, \texttt{collector.cpp} \\
    Bộ máy pipeline       & \texttt{pipeline}  & C++17          & \texttt{DaemonSet} (sidecar)   & \texttt{engine.cpp}, \texttt{libdsa/} \\
    Cầu nối cảnh báo      & \texttt{soc}       & Go             & \texttt{DaemonSet} (sidecar)   & \texttt{alert-bridge/main.go} \\
    Bộ tổng hợp           & \texttt{soc}       & Go             & \texttt{Deployment}            & \texttt{aggregator/correlator.go} \\
    Dịch vụ sự cố         & \texttt{soc}       & Go             & \texttt{Deployment}            & \texttt{incident-service/internal/*} \\
    Console web           & \texttt{soc}       & Python (Flask) & \texttt{Deployment}            & \texttt{web-console/app.py} \\
    Lưu trữ               & \texttt{soc}       & SQL            & \texttt{StatefulSet}           & \texttt{store/migrations/*.sql} \\
    Hạ tầng \& triển khai & \texttt{infra}, \texttt{deploy} & Ansible, YAML & ArgoCD, Tekton & \texttt{ansible/}, \texttt{argocd/} \\
    \bottomrule
  \end{tabularx}
\end{table}

Luồng dữ liệu xuyên suốt các thành phần này như sau: lời gọi \texttt{connect()} của
một pod $\rightarrow$ bộ thu thập eBPF (tầng nhân) $\rightarrow$ pipeline (phát hiện)
$\rightarrow$ cảnh báo qua NATS $\rightarrow$ bộ tổng hợp (khử trùng lặp, tương quan)
$\rightarrow$ dịch vụ sự cố (tạo sự cố, sinh chính sách cách ly) $\rightarrow$ người
vận hành duyệt trên console $\rightarrow$ áp \texttt{NetworkPolicy} vào cụm.
Hình~\ref{fig:arch-overview} cho cái nhìn tổng quát về toàn bộ kiến trúc: mặt phẳng
phát hiện chạy phân tán trên mỗi node (DaemonSet), trung tâm SOC xử lý tập trung, tầng
quan trắc và thực thi chính sách, cùng nền tảng GitOps/CI ở dưới. Các mục tiếp theo
trình bày lần lượt từng mắt xích.

\begin{figure}[H]
  \centering
  \resizebox{\linewidth}{!}{% So do tong quat kien truc he thong K8sSecDaP (phong cach Hinh 34).
% Logo PNG o ../shared/figures/logos/ (graphicspath da tro toi shared/figures).
\begin{tikzpicture}[
  font=\footnotesize, >={Stealth[length=2.2mm]},
  svc/.style ={draw, rounded corners=2pt, fill=white, align=center,
               minimum width=2.9cm, minimum height=1.3cm, inner sep=2pt},
  store/.style={draw, rounded corners=2pt, fill=blue!5, align=center,
               minimum width=2.5cm, minimum height=1.2cm, inner sep=2pt},
  bus/.style ={draw, rounded corners=2pt, fill=green!7, align=center,
               minimum width=1.8cm, minimum height=3.6cm},
  flow/.style ={->, thick, draw=black!75},
  dflow/.style={->, very thick, draw=orange!80!black},
  ctrl/.style ={->, thick, draw=blue!60!black, dashed},
  band/.style ={font=\footnotesize\bfseries, anchor=west, text=black!65},
]
% ---------- Actors (trai) ----------
\node[draw, fill=yellow!12, rounded corners=2pt, align=center, minimum width=2.3cm,
      minimum height=1.0cm] (admin) at (0,6.6) {Quản trị viên\\/ Analyst};
\node[draw, fill=red!8, rounded corners=2pt, align=center, minimum width=2.3cm,
      minimum height=1.1cm] (pods) at (0,2.2) {Pod ứng dụng\\/ Pod tấn công};
\node[svc, fill=cyan!7, minimum width=2.0cm] (ingress) at (3.0,4.4)
  {Cloudflared\\$+$ Traefik\\\scriptsize(ingress)};

% ================= CUM KUBERNETES =================
% --- Bang 1: phat hien (DaemonSet) ---
\node[svc] (collector) at (7.2,8.0) {collector (eBPF)\\
\includegraphics[height=0.46cm]{logos/ebpf.png}};
\node[svc] (pipeline)  at (10.6,8.0) {pipeline\\\scriptsize CMS\,·\,Tarjan\,·\,BFS\,·\,LPM};
\node[svc] (bridge)    at (14.0,8.0) {alert-bridge};

% --- Bang 2: SOC ---
\node[svc, fill=cyan!4] (console) at (14.0,5.7) {web-console\\\scriptsize Flask + HTMX};
\node[bus] (nats) at (7.2,3.7) {NATS\\JetStream\\
\includegraphics[height=0.46cm]{logos/nats.png}};
\node[svc] (agg)  at (10.6,4.9) {aggregator\\\scriptsize tương quan / cross-node};
\node[svc] (inc)  at (14.0,3.5) {incident-service\\\scriptsize sự cố·chính sách·RBAC};
\node[store] (pg) at (17.3,3.5) {PostgreSQL\\
\includegraphics[height=0.46cm]{logos/postgresql.png}};

% --- Bang 3: quan trac + thuc thi ---
\node[svc, minimum width=2.4cm] (prom) at (9.0,1.0) {Prometheus\\
\includegraphics[height=0.42cm]{logos/prometheus.png}};
\node[svc, minimum width=2.4cm] (graf) at (12.0,1.0) {Grafana\\
\includegraphics[height=0.42cm]{logos/grafana.png}};
\node[svc, minimum width=2.4cm] (k8s)  at (17.3,1.0) {K8s API + Calico\\
\includegraphics[height=0.42cm]{logos/calico.png}};
\node[draw, fill=red!7, rounded corners=2pt, align=center, minimum width=2.2cm,
      minimum height=0.9cm] (quar) at (17.3,-0.8) {Pod bị cách ly\\\scriptsize (egress-deny)};

% --- Hop cum + nhan ---
\node[draw, dashed, very thick, draw=black!55, rounded corners=6pt, inner sep=16pt,
      fit=(collector)(pipeline)(bridge)(console)(nats)(agg)(inc)(pg)(prom)(graf)(k8s)(quar)] (zk8s) {};
\node[font=\small\bfseries, anchor=north east] at ([xshift=-4pt,yshift=-2pt]zk8s.north east)
  {Cụm Kubernetes \;
\includegraphics[height=0.42cm]{logos/kubernetes.png}};
\node[band] at ([yshift=10pt]collector.north west) {Mặt phẳng phát hiện — DaemonSet mỗi node};
\node[band] at ([xshift=-6pt,yshift=20pt]nats.north west) {Trung tâm vận hành an ninh (SOC)};
\node[band] at ([xshift=-6pt,yshift=14pt]prom.north west) {Quan trắc \& Thực thi chính sách};

% ---------- Nen tang GitOps/CI (duoi) ----------
\node[svc, minimum width=2.3cm] (argo)   at (7.2,-3.6) {ArgoCD\\
\includegraphics[height=0.42cm]{logos/argo.png}};
\node[svc, minimum width=2.3cm] (tekton) at (10.6,-3.6) {Tekton (CI)\\
\includegraphics[height=0.42cm]{logos/tekton.png}};
\node[store, minimum width=2.3cm] (harbor) at (14.0,-3.6) {Harbor\\
\includegraphics[height=0.42cm]{logos/harbor.png}};
\node[draw, dashed, draw=black!55, rounded corners=6pt, inner sep=12pt,
      fit=(argo)(tekton)(harbor)] (zplat) {};
\node[font=\footnotesize\bfseries, text=black!65, anchor=north west]
  at ([xshift=4pt,yshift=-3pt]zplat.south west) {Nền tảng GitOps / CI};

% ================= LUONG =================
\draw[flow] (admin) -| (ingress.north);
\draw[flow] (ingress) |- (console.west);
\draw[<->, thick, draw=black!75] (console.south) -- node[right,font=\scriptsize]{REST} (inc.north);
% data plane (cam)
\draw[dflow] (pods.east) -- ++(0.5,0) |- node[pos=0.78,above,font=\scriptsize]{\texttt{connect()}} (collector.west);
\draw[dflow] (collector) -- (pipeline);
\draw[dflow] (pipeline) -- (bridge);
\draw[dflow] (bridge.south) -- (14,6.55) -- node[above,font=\scriptsize]{\texttt{zt.alerts.raw}} (7.2,6.55) -- (nats.north);
% SOC
\draw[dflow] (nats) -- node[above,font=\scriptsize]{alerts} (agg);
\draw[flow] (agg) -- node[above,font=\scriptsize]{enriched} (inc);
\draw[flow] (inc) -- (pg);
% thuc thi (control — xanh)
\draw[ctrl] (inc) -- node[right,font=\scriptsize]{áp NetworkPolicy} (k8s);
\draw[ctrl] (k8s) -- (quar);
% quan trac
\draw[flow] (pipeline.south) -- (prom.north);
\draw[flow] (prom) -- (graf);
% GitOps
\draw[ctrl] (argo.north) -- node[left,font=\scriptsize]{sync} (zk8s.south-|argo.north);
\draw[flow] (tekton) -- node[below,font=\scriptsize]{push} (harbor);
\draw[ctrl] (harbor.north) -- node[right,font=\scriptsize]{pull image} (zk8s.south-|harbor.north);
\end{tikzpicture}
}
  \caption{Sơ đồ tổng quát kiến trúc hệ thống.}
  \label{fig:arch-overview}
\end{figure}

\section{Môi trường và công cụ triển khai}
\label{sec:impl-env}

Hệ thống được hiện thực và kiểm thử trên một cụm Kubernetes thật gồm ba node (một
node control-plane và hai node worker), chạy hệ điều hành Ubuntu~24.04 với nhân Linux
phiên bản~6.8 có sẵn BTF, runtime container \texttt{containerd} và mạng pod do CNI
Calico đảm nhiệm (có hỗ trợ \texttt{NetworkPolicy}). Mỗi tầng được hiện thực bằng ngôn
ngữ phù hợp với ràng buộc của tầng đó: chương trình tầng nhân viết bằng C và biên dịch
sang đối tượng eBPF bằng \texttt{clang/LLVM}; bộ đọc và bộ máy pipeline viết bằng
C++17 để kiểm soát chặt bộ nhớ và hiệu năng; các dịch vụ SOC viết bằng Go; console
web viết bằng Python (Flask). Mỗi thành phần được đóng gói thành ảnh OCI riêng và nạp
vào cụm qua đối tượng Kubernetes tương ứng (\texttt{DaemonSet}, \texttt{Deployment},
\texttt{StatefulSet}). Việc tách bạch ngôn ngữ và đóng gói theo tầng giúp cô lập phần
chạy đặc quyền ở tầng nhân khỏi phần logic nghiệp vụ.

Tôi dựng cụm tự động bằng Ansible (kho \texttt{infra}): mỗi node được khởi tạo bằng
\texttt{kubeadm}, cài \texttt{containerd}, rồi CNI Calico được áp qua role
\texttt{cni\_calico}. Sau khi dựng xong, tôi kiểm chứng môi trường bằng
\texttt{kubectl get nodes} (Mã nguồn~\ref{lst:get-nodes}): cả ba node ở trạng thái
\texttt{Ready}, cùng nhân Linux 6.8 và runtime \texttt{containerd} --- xác nhận điều
kiện chạy eBPF CO-RE (nhân $\geq 5.8$, có BTF) đã được thỏa mãn đồng nhất trên cụm.

\begin{lstlisting}[language={},caption={Kiểm chứng cụm sau khi dựng bằng Ansible (kết quả \texttt{kubectl} rút gọn).},label={lst:get-nodes}]
$ kubectl get nodes -o wide
NAME               STATUS  ROLES          VERSION  KERNEL-VERSION    CONTAINER-RUNTIME
userver-master     Ready   control-plane  v1.32.13  6.8.0-generic     containerd://2.2.4
userver-worker-1   Ready   <none>         v1.32.13  6.8.0-generic     containerd://2.2.4
userver-worker-home   Ready   <none>         v1.32.13  6.8.0-generic     containerd://2.2.4
\end{lstlisting}

\section{Hạ tầng nền tảng và triển khai (giới thiệu)}
\label{sec:impl-platform}

Phần này được trình bày vắn tắt vì không thuộc trọng tâm đề tài, nhưng vẫn nêu đủ các
thành phần để bức tranh triển khai hoàn chỉnh. Toàn bộ trạng thái mong muốn của hệ
thống được quản lý theo mô hình GitOps: manifest Kubernetes lưu trong kho Git và một
ứng dụng ArgoCD (\texttt{zt-soc-core}) liên tục đồng bộ trạng thái cụm về đúng trạng
thái khai báo, với hai tùy chọn \textit{prune} và \textit{self-heal} bật để tự dọn tài
nguyên thừa và tự khôi phục khi có sai lệch. Quy trình tích hợp liên tục do pipeline
\emph{Tekton} chạy ngay trong cụm đảm nhiệm: clone mã nguồn, dựng ảnh bằng Kaniko,
đẩy lên kho ảnh riêng (\emph{Harbor}), rồi cập nhật thẻ ảnh trong manifest để ArgoCD
triển khai. Hạ tầng node (cài đặt \texttt{kubeadm}, \texttt{containerd}, CNI Calico)
được tự động hóa bằng Ansible trong kho \texttt{infra}.

Cụ thể, gốc của cây GitOps là một \texttt{Application} kiểu ``app-of-apps'': nó trỏ vào
thư mục \texttt{platform/argocd/apps} và sinh ra một \texttt{Application} con cho mỗi
tệp, trong đó thứ tự triển khai được kiểm soát bằng \emph{sync-wave} để CNI luôn tồn
tại trước khi pod SOC xin địa chỉ pod-network (Mã nguồn~\ref{lst:argocd}). Nhờ
\textit{selfHeal}, mọi sai lệch thủ công đều bị kéo về đúng trạng thái khai báo trong
Git.

\begin{lstlisting}[language={},caption={Gốc GitOps ``app-of-apps'' và thứ tự sync-wave của các ứng dụng con (\texttt{platform/argocd/}).},label={lst:argocd}]
# app-of-apps.yaml: goc, sinh 1 Application cho moi tep trong apps/
kind: Application
metadata: { name: platform-apps, namespace: argocd }
spec:
  source: { repoURL: ...K8sSecDaP-deploy.git, path: platform/argocd/apps }
  syncPolicy: { automated: { prune: true, selfHeal: true } }
---
# Thu tu trien khai bang sync-wave (CNI truoc, SOC sau, lop du lieu cuoi):
#   zt-cni-calico   sync-wave: "-10"   (prune/selfHeal=false — KHONG bao gio prune CNI)
#   zt-soc-core     sync-wave: "0"     -> manifests/ cua K8sSecDaP-soc
#   zt-soc-data     sync-wave: "10"    -> data/ (MinIO, cronjob bao cao tuan)
\end{lstlisting}

Quy trình CI là một Tekton \texttt{Pipeline} tên \texttt{build-and-deploy} gồm bốn
\texttt{Task} tuần tự: \texttt{git-clone} $\rightarrow$ \texttt{kaniko-build} (dựng ảnh
OCI không cần Docker daemon) $\rightarrow$ đẩy ảnh lên kho riêng $\rightarrow$
\texttt{update-manifest} (vá thẻ ảnh vào manifest GitOps, commit để ArgoCD đồng bộ).
Như vậy một thay đổi mã nguồn được đưa ra cụm hoàn toàn tự động và có thể truy vết qua
Git.

Để minh chứng các thành phần này thực sự đang chạy, tôi truy vấn trạng thái cụm
(Mã nguồn~\ref{lst:argocd-proof}): các \texttt{Application} của ArgoCD đều
\texttt{Synced/Healthy}, và \texttt{DaemonSet} thu thập \texttt{zt-pipeline-ebpf} hiện
diện trên các node worker.

\begin{lstlisting}[language={},caption={Minh chứng triển khai: ArgoCD đồng bộ và DaemonSet thu thập đang chạy (\texttt{kubectl} rút gọn).},label={lst:argocd-proof}]
$ kubectl get applications -n argocd
NAME            SYNC STATUS   HEALTH STATUS
platform-apps   Synced        Healthy
zt-cni-calico   Synced        Healthy
zt-soc-core     Synced        Healthy
zt-soc-data     Synced        Healthy

$ kubectl get daemonset -n zt-mapper
NAME               DESIRED   CURRENT   READY   NODE SELECTOR
zt-pipeline-ebpf   2         2         2       <none>
\end{lstlisting}

Ở tầng truyền thông, các dịch vụ trao đổi cảnh báo qua trục thông điệp NATS theo các
chủ đề có cấu trúc: \texttt{zt.alerts.raw} (đầu ra pipeline mỗi node) và
\texttt{zt.alerts.enriched} (sau khi tổng hợp). Ở tầng lưu trữ, PostgreSQL chạy dưới
dạng \texttt{StatefulSet} có ổ đĩa bền vững để giữ sự cố, cảnh báo và nhật ký kiểm
toán. Hệ thống quan trắc Prometheus/Grafana cung cấp số liệu và dashboard cho đội vận
hành. Hình~\ref{fig:deploy-topology} tóm tắt cách các thành phần phân bố trên cụm ba
node cùng vòng đồng bộ GitOps.

\begin{figure}[H]
  \centering
  \resizebox{0.95\linewidth}{!}{% Sơ đồ triển khai (topology) cụm K8s 3 node + GitOps. Dùng ở Chương 4.
\begin{tikzpicture}[
    node distance=0.5cm,
    every node/.style={font=\footnotesize},
    node3/.style={draw, rounded corners=2pt, fill=blue!6, align=center,
                  minimum width=3.2cm, minimum height=1.7cm},
    pod/.style={draw, rounded corners=1pt, fill=red!12, align=center,
                font=\scriptsize, minimum width=2.6cm, minimum height=0.55cm},
    svc/.style={draw, rounded corners=1pt, fill=green!12, align=center,
                font=\scriptsize, minimum width=2.4cm, minimum height=0.55cm},
    flow/.style={-{Stealth[length=2mm]}, thick},
  ]
  \node[node3, label=above:\textbf{master (control-plane)}] (m) {};
  \node[pod] at ([yshift=0.35cm]m.center) (mc) {collector (DaemonSet)};
  \node[svc] at ([yshift=-0.35cm]m.center) (ms) {NATS · Postgres · SOC};

  \node[node3, right=1.0cm of m, label=above:\textbf{worker-1}] (w1) {};
  \node[pod] at ([yshift=0.35cm]w1.center) (w1c) {collector (DaemonSet)};
  \node[svc] at ([yshift=-0.35cm]w1.center) (w1s) {Pod mục tiêu (web/api)};

  \node[node3, right=1.0cm of w1, label=above:\textbf{worker-home}] (w2) {};
  \node[pod] at ([yshift=0.35cm]w2.center) (w2c) {collector (DaemonSet)};
  \node[svc] at ([yshift=-0.35cm]w2.center) (w2s) {Pod ứng dụng};

  \node[draw, rounded corners, fill=orange!14, align=center, below=1.3cm of w1,
        minimum width=3.4cm] (argo) {ArgoCD (GitOps)\\Helm · Calico CNI};

  \draw[flow] (w1c) -- (ms);
  \draw[flow] (w2c) to[bend left=10] (ms);
  \draw[flow] (argo) -- node[right,font=\scriptsize]{đồng bộ manifest} (w1);
  \draw[flow] (argo) to[bend right=12] (m);
\end{tikzpicture}
}
  \caption{Sơ đồ triển khai trên cụm Kubernetes ba node.}
  \label{fig:deploy-topology}
\end{figure}

\section{Bộ thu thập eBPF}
\label{sec:impl-collector}

Bộ thu thập là mắt xích đầu tiên và là thành phần đặc quyền duy nhất chạm vào tầng
nhân, nên được trình bày kỹ nhất. Về kiến trúc, nó gồm hai phần (đã mô tả ở
Hình~\ref{fig:comp-collect-pipe}, Chương~\ref{ch:design}): \emph{chương trình tầng
nhân} gắn \texttt{kprobe}/\texttt{kretprobe} vào hàm \texttt{tcp\_v4\_connect} và lọc
sự kiện ngay trong nhân; và \emph{chương trình tầng người dùng} nạp, gắn chương trình
eBPF, đọc sự kiện từ bộ đệm vòng rồi xuất ra dòng JSON.

\subsection{Chương trình tầng nhân}
\label{subsec:impl-kern}

Chương trình tầng nhân khai báo ba BPF map, mỗi map giữ một vai trò rõ rệt:

\begin{itemize}[leftmargin=1.4em]
  \item \texttt{cidr\_map} (\texttt{LPM\_TRIE}, 1024 mục): ánh xạ tiền tố CIDR sang
        nhãn (1=pod, 2=service, 3=node, 4=RFC1918, 5=loopback) để phân loại địa chỉ
        \emph{ngay trong nhân}. Bắt buộc cờ \texttt{BPF\_F\_NO\_PREALLOC}.
  \item \texttt{events} (\texttt{RINGBUF}, 256~KiB): kênh truyền sự kiện một chiều
        từ nhân lên không gian người dùng, thay cho perf buffer cũ, cho thông lượng
        cao hơn và không mất thứ tự.
  \item \texttt{sock\_store} (\texttt{HASH}, khóa \texttt{pid\_tgid}): lưu tạm con
        trỏ \texttt{struct sock} giữa lúc \texttt{kprobe} vào và \texttt{kretprobe}
        ra, vì thông tin địa chỉ đích chỉ đầy đủ \emph{sau} khi kết nối được thiết lập.
\end{itemize}

Cơ chế hai pha (entry $\rightarrow$ return) là điểm tinh tế nhất: \texttt{kprobe} lúc
vào chỉ lưu con trỏ socket theo \texttt{pid\_tgid}; \texttt{kretprobe} lúc ra mới đọc
địa chỉ nguồn/đích bằng CO-RE (\texttt{BPF\_CORE\_READ}), \emph{lọc nguồn} (chỉ phát
sự kiện có nhãn nguồn = 1, tức xuất phát từ mạng pod~--~đúng mô hình mối đe dọa
``pod bị chiếm''), rồi đặt sự kiện vào bộ đệm vòng. Mã
nguồn~\ref{lst:bpf-kretprobe} trích phần lõi của \texttt{kretprobe}.

\begin{lstlisting}[language=C,caption={Trích \texttt{kretprobe}: đọc socket bằng CO-RE, lọc nguồn theo LPM, đẩy sự kiện qua ring buffer (\texttt{tcp\_connect.bpf.c}).},label={lst:bpf-kretprobe}]
SEC("kretprobe/tcp_v4_connect")
int BPF_KRETPROBE(tcp_v4_connect_return, int ret)
{
    if (ret != 0 && ret != -115)   // chi xu ly connect() thanh cong / EINPROGRESS
        return 0;
    __u64 pid_tgid = bpf_get_current_pid_tgid();
    struct sock **skp = bpf_map_lookup_elem(&sock_store, &pid_tgid);
    if (!skp) return 0;
    struct sock *sk = *skp;
    bpf_map_delete_elem(&sock_store, &pid_tgid);

    __u32 dst_ip = 0, src_ip = 0; __u16 dst_port = 0;
    BPF_CORE_READ_INTO(&dst_ip,  sk, __sk_common.skc_daddr);
    BPF_CORE_READ_INTO(&dst_port, sk, __sk_common.skc_dport);
    BPF_CORE_READ_INTO(&src_ip,  sk, __sk_common.skc_rcv_saddr);
    dst_port = __builtin_bswap16(dst_port);

    // Loc nguon: chi phat su kien neu nguon thuoc mang pod (nhan == 1).
    struct lpm_key src_key = { .prefixlen = 32, .addr = src_ip };
    __u32 *src_label = bpf_map_lookup_elem(&cidr_map, &src_key);
    if (!src_label || *src_label != 1)
        return 0;

    struct network_event_t *evt = bpf_ringbuf_reserve(&events, sizeof(*evt), 0);
    if (!evt) return 0;
    evt->src_ip = src_ip; evt->dst_ip = dst_ip; evt->dst_port = dst_port;
    evt->pid = pid_tgid >> 32;
    evt->timestamp_ns = bpf_ktime_get_ns();
    bpf_ringbuf_submit(evt, 0);
    return 0;
}
\end{lstlisting}

Việc lọc nguồn ngay trong nhân là một quyết định thiết kế quan trọng: nó loại bỏ phần
lớn lưu lượng không liên quan (kết nối do chính node phát đi, lưu lượng host) \emph{tại
nguồn}, giảm số sự kiện phải truyền lên không gian người dùng và~--~như phân tích ở
Chương~\ref{ch:eval}~--~giảm đáng kể tỷ lệ báo động giả. Đường đi đầy đủ của một gói
tin từ lời gọi \texttt{connect()} tới khi thành sự kiện ở không gian người dùng được
mô tả ở Hình~\ref{fig:packet-event-path}.

\begin{figure}[H]
  \centering
  \resizebox{0.9\linewidth}{!}{% Duong di chi tiet: connect() -> kprobe/kretprobe -> loc LPM -> RINGBUF -> collector -> pipeline. Chuong 2.
\begin{tikzpicture}[font=\footnotesize, node distance=0.62cm and 0.9cm,
    pod/.style ={draw, rounded corners=2pt, fill=blue!10, align=center, minimum height=0.8cm, minimum width=4.4cm},
    ker/.style ={draw, rounded corners=2pt, fill=green!12, align=center, minimum height=0.8cm, minimum width=5.6cm},
    dec/.style ={draw, diamond, aspect=2.6, fill=yellow!20, align=center, inner sep=1pt},
    usr/.style ={draw, rounded corners=2pt, fill=blue!8, align=center, minimum height=0.8cm, minimum width=5.6cm},
    map/.style ={draw, rounded corners=2pt, fill=orange!16, align=center, minimum height=0.7cm, font=\scriptsize},
    flow/.style={-{Stealth[length=2mm]}, thick},
  ]
  \node[pod] (a) {\textbf{Pod:} \texttt{connect(dst\_ip:dst\_port)}};
  \node[ker, below=of a] (b) {\textbf{kprobe} \texttt{tcp\_v4\_connect\_entry}:\\lưu con trỏ \texttt{sock} theo \texttt{pid\_tgid}};
  \node[ker, below=of b] (c) {\textbf{kretprobe} \texttt{tcp\_v4\_connect\_return}:\\đọc \texttt{sock} $\rightarrow$ trích \texttt{src\_ip, dst\_ip, dst\_port}};
  \node[ker, below=of c] (d) {Tra \texttt{cidr\_map} (LPM-Trie)\\$\rightarrow$ \texttt{src\_label}, \texttt{dst\_label}};
  \node[dec, below=of d] (e) {\texttt{src\_label}\\$=1$ (pod)?};
  \node[ker, below=1.0cm of e] (f) {Cấp ô \texttt{events} (RINGBUF), điền\\\texttt{network\_event\_t}\{src,dst,port,pid,ts\} $\rightarrow$ submit};
  \node[usr, below=of f] (g) {\textbf{Collector} (libbpf, \texttt{ring\_buffer\_\_poll})\\$\rightarrow$ dòng JSON $\rightarrow$ FIFO $\rightarrow$ \textbf{pipeline}};
  \node[draw, rounded corners=2pt, fill=gray!12, right=of e, align=center, minimum height=0.7cm] (drop) {bỏ qua\\(không phải pod)};

  % maps on the side
  \node[map, right=0.7cm of b] (m1) {\texttt{sock\_store}\\HASH};
  \node[map, right=0.7cm of d] (m2) {\texttt{cidr\_map}\\LPM\_TRIE};
  \node[map, right=0.7cm of f] (m3) {\texttt{events}\\RINGBUF};

  \draw[flow] (a)--(b); \draw[flow] (b)--(c); \draw[flow] (c)--(d); \draw[flow] (d)--(e);
  \draw[flow] (e)-- node[right,font=\scriptsize]{đúng} (f);
  \draw[flow] (e)-- node[above,font=\scriptsize]{sai} (drop);
  \draw[flow] (f)--(g);
  \draw[{Stealth}-{Stealth}, gray] (b.east)--(m1.west);
  \draw[{Stealth}-{Stealth}, gray] (d.east)--(m2.west);
  \draw[{Stealth}-{Stealth}, gray] (f.east)--(m3.west);
\end{tikzpicture}
}
  \caption{Đường đi của một sự kiện kết nối qua eBPF.}
  \label{fig:packet-event-path}
\end{figure}

\subsection{Chương trình tầng người dùng}
\label{subsec:impl-user}

Phần tầng người dùng (\texttt{collector.cpp}) dùng \texttt{libbpf} để: mở tệp đối
tượng \texttt{tcp\_connect.bpf.o}, nạp vào nhân, gắn \texttt{kprobe}/\texttt{kretprobe},
\emph{nạp bảng CIDR} vào \texttt{cidr\_map}, rồi vào vòng lặp đọc bộ đệm vòng. Bảng
CIDR (Mã nguồn~\ref{lst:user-cidr}) chính là tri thức để chương trình tầng nhân phân
loại địa chỉ; mỗi mục là cặp (tiền tố, nhãn) với địa chỉ ở thứ tự byte mạng.

\begin{lstlisting}[language=C++,caption={Nạp bảng CIDR vào \texttt{LPM\_TRIE} từ tầng người dùng (\texttt{collector.cpp}).},label={lst:user-cidr}]
struct { struct lpm_key key; uint32_t label; const char *desc; } entries[] = {
    { {16, htonl(0x0AF40000)}, 1, "pod_network" },     // 10.244.0.0/16
    { {12, htonl(0x0A600000)}, 2, "service_network" },  // 10.96.0.0/12
    { {16, htonl(0xC0A80000)}, 3, "node_network" },     // 192.168.0.0/16
    { {12, htonl(0xAC100000)}, 4, "private_172" },      // 172.16.0.0/12
    { {8,  htonl(0x7F000000)}, 5, "loopback" },         // 127.0.0.0/8
};
for (const auto &e : entries)
    bpf_map_update_elem(map_fd, &e.key, &e.label, BPF_ANY);
\end{lstlisting}

Mỗi sự kiện đọc được từ bộ đệm vòng được chuyển con trỏ qua hàm callback
\texttt{handle\_event}, ở đó địa chỉ nhị phân được đổi sang chuỗi chấm-thập phân và in
ra \texttt{stdout} dưới dạng một dòng JSON (ví dụ
\texttt{\{"src\_ip":"10.244.1.5", "dst\_port":8080, \dots\}}). Vòng
lặp chính gọi \texttt{ring\_buffer\_\_poll} với chu kỳ 100~ms và thoát sạch khi nhận
tín hiệu \texttt{SIGINT}/\texttt{SIGTERM}. Việc xuất JSON theo dòng giúp giữ giao diện
\emph{đơn giản và thay thế được}: pipeline chỉ cần đọc \texttt{stdin}, không phụ thuộc
vào thư viện eBPF.

\subsection{Đóng gói và triển khai}
\label{subsec:impl-collector-deploy}

Để chạy được trên cụm có nhân không đồng nhất, ảnh Docker của bộ thu thập sinh tệp
\texttt{vmlinux.h} từ thông tin BTF của nhân ngay trong quá trình xây dựng, hiện thực
hóa cơ chế CO-RE đã trình bày ở Chương~\ref{ch:foundations}. Bộ thu thập được triển
khai dưới dạng \texttt{DaemonSet} \emph{đặc quyền} để hiện diện trên mọi node. Trong
cùng pod, ba container chạy cạnh nhau và nối với nhau qua một FIFO trên
\texttt{emptyDir}: \texttt{collector} ghi JSON vào FIFO, \texttt{pipeline} đọc FIFO
qua \texttt{stdin}, còn \texttt{alert-bridge} đẩy cảnh báo lên NATS. Mã
nguồn~\ref{lst:collector-ds} trích phần khai báo container thu thập, bao gồm cơ chế
tự khởi động lại và \texttt{securityContext} đặc quyền.

\begin{lstlisting}[language={},caption={Trích \texttt{DaemonSet} thu thập: tạo FIFO, vòng tự khởi động lại, quyền đặc quyền nạp eBPF (\texttt{23-pipeline-ebpf.yaml}).},label={lst:collector-ds}]
- name: collector
  image: johnnaeder/ebpf-portscan-k8s-check:collector-ebpf-s3
  command: ["/bin/sh", "-c"]
  args:
    - |
      FIFO=/stream/events.pipe
      [ -p "$FIFO" ] || mkfifo -m 0600 "$FIFO"
      while true; do
        /app/collector > "$FIFO" || echo "[ebpf] collector exited $?" >&2
        sleep 3
      done
  volumeMounts:
    - { name: stream, mountPath: /stream }
    - { name: btf,    mountPath: /sys/kernel/btf, readOnly: true }
  securityContext:
    privileged: true        # can de nap + attach BPF, doc kernel BTF
\end{lstlisting}

Vì \texttt{collector} chạy đặc quyền và nạp mã eBPF vào nhân, đây là thành phần có quyền
cao hơn hẳn ứng dụng thường, nên tôi đặt nó vào một ranh giới tin cậy riêng và làm cứng
tách biệt; phần phân tích mô hình tin cậy và các biện pháp làm cứng được trình bày cùng
với các giới hạn an ninh ở mục~\ref{sec:eval-trust} (Chương~\ref{ch:eval}).

\subsection{Kiểm chứng eBPF hoạt động ở tầng nhân}
\label{subsec:impl-collector-verify}

Vì can thiệp tầng nhân là trọng tâm của đồ án, cần một cách \emph{quan sát trực tiếp}
chứng minh chương trình eBPF đã thực sự được nạp, kiểm chứng và gắn vào ngăn xếp mạng
của nhân~--~chứ không chỉ là một tiến trình người dùng thông thường. Việc kiểm chứng dựa
trên \texttt{bpftool} (truy vấn trạng thái BPF của nhân) và giao diện tracing của nhân,
gồm ba bước minh chứng độc lập (Mã nguồn~\ref{lst:bpf-verify}).

\emph{Thứ nhất}, \texttt{bpftool prog show} liệt kê chương trình đang nạp: kiểu
\texttt{kprobe} gắn vào hàm \texttt{tcp\_v4\_connect}~--~đúng điểm móc đã chọn ở
Chương~\ref{ch:foundations}~--~kèm số lệnh đã được trình kiểm chứng (verifier) duyệt và
đã biên dịch JIT, khẳng định mã chạy \emph{trong} nhân chứ không phải trong không gian
người dùng. \emph{Thứ hai}, \texttt{bpftool map show} liệt kê đúng ba bản đồ BPF mà bộ
thu thập sử dụng: \texttt{cidr\_map} kiểu \texttt{lpm\_trie} (lọc dải pod-network),
\texttt{events} kiểu \texttt{ringbuf} (kênh chuyển sự kiện lên không gian người dùng) và
\texttt{sock\_store} kiểu \texttt{hash} (ghép pha vào/ra của \texttt{connect}).
\emph{Thứ ba}, đổ nội dung \texttt{cidr\_map} xác nhận các CIDR pod-network đã được nạp,
và đọc luồng sự kiện cho thấy mỗi lời gọi \texttt{connect()} trong cụm sinh đúng một bản
ghi sự kiện~--~tức là cơ chế tracing đang chạy theo thời gian thực.

\begin{lstlisting}[language={},caption={Kiểm chứng chương trình và bản đồ eBPF ở tầng nhân (chạy \texttt{bpftool} trên node).},label={lst:bpf-verify}]
# (1) Hai chuong trinh eBPF da nap: kprobe + kretprobe gan tcp_v4_connect, da JIT
$ sudo bpftool prog show
97069: kprobe  name tcp_v4_connect_entry   gpl  run_cnt 101167
        xlated 112B  jited 72B   memlock 4096B  map_ids 19041              btf_id 3973
97070: kprobe  name tcp_v4_connect_return  gpl  run_cnt 101167
        xlated 760B  jited 412B  memlock 4096B  map_ids 19041,19039,19040  btf_id 3973

# (2) Ba ban do BPF cua bo thu thap
$ sudo bpftool map show
19039: lpm_trie  name cidr_map    flags 0x1  key 8B  value 4B  max_entries 1024
19040: ringbuf   name events      flags 0x0                    max_entries 262144
19041: hash      name sock_store  flags 0x0  key 8B  value 8B  max_entries 4096

# (3) Cac CIDR pod-network da nap vao bo loc LPM
$ sudo bpftool map dump name cidr_map | head
[{ "key": { "prefixlen": 12, "addr": 24586 }, "value": 2 },
 { "key": { "prefixlen": 16, "addr": 62474 }, "value": 1 }]

# (4) Tracing thoi gian thuc: moi connect() sinh mot su kien JSON tren collector stdout
$ kubectl -n zt-mapper logs zt-pipeline-ebpf-tx94l -c collector --tail=1
{"src":"10.244.204.189","dst":"10.244.204.161","dport":80,"src_label":1}
\end{lstlisting}

Bốn đầu ra trên là minh chứng đầy đủ cho chuỗi ``nạp~$\rightarrow$~kiểm chứng~$\rightarrow$
JIT~$\rightarrow$~gắn móc~$\rightarrow$~bắt sự kiện'' diễn ra ở tầng nhân.
Hình~\ref{fig:ebpf-bpftool} là ảnh chụp đầu ra \texttt{bpftool} thật trên node
\texttt{userver-worker-1}: hai chương trình \texttt{kprobe}/\texttt{kretprobe} gắn vào
\texttt{tcp\_v4\_connect} (đã \texttt{xlated}/\texttt{jited}, số lần chạy \texttt{run\_cnt}
hơn $10^5$) và đúng ba bản đồ \texttt{cidr\_map} (\texttt{lpm\_trie}), \texttt{events}
(\texttt{ringbuf}), \texttt{sock\_store} (\texttt{hash}). Quy trình chạy lại đầy đủ được
tổng hợp trong Phụ lục~\ref{app:demo}.

\begin{figure}[H]
  \centering
  \IfFileExists{assets/images/ebpf-bpftool.png}{%
    \includegraphics[width=0.92\linewidth]{ebpf-bpftool}}{%
    \fbox{\parbox[c][3.2cm][c]{0.8\linewidth}{\centering\itshape\small
    [Ảnh chụp đầu ra \texttt{bpftool} trên node --- lưu vào
    \texttt{stage2/assets/images/ebpf-bpftool.png}]}}}
  \caption{Minh chứng eBPF ở tầng nhân qua \texttt{bpftool} trên node.}
  \label{fig:ebpf-bpftool}
\end{figure}

\section{Bộ máy pipeline}
\label{sec:impl-pipeline}

Bộ máy pipeline (\texttt{engine.cpp} + thư viện \texttt{libdsa}) là nơi đặt các thuật
toán phát hiện. Kiến trúc của nó (Hình~\ref{fig:comp-collect-pipe}, phải) gồm một bộ
\emph{sổ đăng ký} (Registry) dựng các thuật toán theo tên trong cấu hình, và một vòng
lặp xử lý sự kiện đọc từ \texttt{stdin}.

\subsection{Cấu hình theo dữ liệu}
\label{subsec:impl-pipeline-config}

Toàn bộ hành vi phát hiện~--~kích thước Count-Min Sketch, ngưỡng, độ dài cửa sổ, danh
sách dải mạng phân loại, luật mức độ nghiêm trọng~--~đều khai báo trong một tệp YAML
nạp vào pod qua \texttt{ConfigMap} (Mã nguồn~\ref{lst:pipeline-cfg}), nên có thể điều
chỉnh mà không cần biên dịch lại. Đây cũng là nơi định nghĩa ánh xạ CIDR $\rightarrow$
nhãn mà pipeline dùng cho phân loại LPM ở tầng người dùng (song song với bảng trong
nhân của bộ thu thập).

\begin{lstlisting}[language={},caption={Trích cấu hình pipeline: thuật toán, ngưỡng, dải mạng, cửa sổ phân tích (\texttt{pipeline.yaml}).},label={lst:pipeline-cfg}]
algorithms:
  frequency_estimator: { name: count_min_sketch, params: { width: 2048, depth: 5, threshold: 40 } }
  scc_detector:        { name: tarjan }
  reachability:        { name: bfs, params: { max_hops: 3 } }
  ip_classifier:       { name: lpm_trie }
networks:
  - { cidr: "10.244.0.0/16", label: 1, name: pod_network }
  - { cidr: "10.96.0.0/12",  label: 2, name: service_network }
  - { cidr: "0.0.0.0/0",     label: 0, name: external }
analysis: { window_seconds: 60, scc_interval_seconds: 10 }
severity_rules:
  - { type: port_scan,    severity: medium }
  - { type: blast_radius, severity: high }
\end{lstlisting}

\subsection{Vòng lặp xử lý sự kiện}
\label{subsec:impl-pipeline-loop}

Với mỗi sự kiện, vòng lặp thực hiện một chuỗi bước cố định (Giải
thuật~\ref{alg:pipeline-loop}): kiểm tra và đặt lại cửa sổ trượt; phân loại địa chỉ
bằng LPM; cập nhật bộ đếm Count-Min Sketch \emph{theo địa chỉ nguồn} và ước lượng tần
suất; nếu vượt ngưỡng thì phát cảnh báo \texttt{port\_scan} (mỗi nguồn chỉ một lần
trong cửa sổ) và tính \emph{bán kính ảnh hưởng} bằng BFS $k$~bước để phát cảnh báo
\texttt{blast\_radius}; bồi đắp đồ thị truyền thông; và định kỳ chạy Tarjan để phát
hiện thành phần liên thông mạnh (\texttt{scc\_anomaly}). Lưu đồ tương ứng ở
Hình~\ref{fig:detect-flow-impl}.

\begin{lstlisting}[style=pseudo,caption={Vòng lặp xử lý sự kiện của pipeline (rút gọn từ \texttt{engine.cpp}).},label={alg:pipeline-loop}]
for each event e read from stdin:
    if e.ts - window_start > window_seconds: reset CMS; clear alerted; window_start = e.ts
    src_label = lpm.lookup(e.src_ip)            // phan loai dia chi nguon
    cms.record(e.src_ip); freq = cms.estimate(e.src_ip)
    if freq > threshold and e.src_ip not in alerted:
        alerted.add(e.src_ip); emit Alert("port_scan", e.src_ip, freq)
        if graph.has(e.src_ip):
            blast = bfs.reachable_within(graph, e.src_ip, max_hops)
            emit Alert("blast_radius", e.src_ip, blast)
    graph.add_edge(e.src_ip, e.dst_ip)          // boi dap do thi truyen thong
    if total_events % (scc_interval*10) == 0:
        sccs = tarjan.find_sccs(graph)
        if new_nontrivial_scc(sccs): emit Alert("scc_anomaly", scc)
\end{lstlisting}

\begin{figure}[H]
  \centering
  \resizebox{0.62\linewidth}{!}{% Luu do xu ly mot su kien ket noi qua cac giai thuat. Dung o Chuong 2.
\begin{tikzpicture}[font=\footnotesize, node distance=0.75cm and 1.4cm,
    start/.style={draw, rounded corners=9pt, fill=gray!12, align=center,
                  minimum height=0.8cm, minimum width=3cm},
    proc/.style ={draw, fill=blue!10, align=center, minimum height=0.8cm,
                  minimum width=3.4cm},
    dec/.style  ={draw, diamond, aspect=2.2, fill=yellow!20, align=center,
                  inner sep=1pt},
    alert/.style={draw, fill=orange!20, align=center, minimum height=0.8cm,
                  minimum width=3.6cm},
    flow/.style ={-{Stealth[length=2mm]}, thick},
  ]
  \node[start] (e) {Sự kiện \texttt{connect()}};
  \node[proc, below=of e] (lpm) {Phân loại IP (LPM Trie)};
  \node[proc, below=of lpm] (cms) {CMS: \texttt{record(src\_ip)}};
  \node[dec, below=of cms] (chk) {ước lượng\\$> \theta$ ?};
  \node[alert, right=of chk] (al) {Cảnh báo \texttt{port\_scan}\\$+$ bán kính ảnh hưởng (BFS)};
  \node[proc, below=1.2cm of chk] (graph) {Cập nhật đồ thị truyền thông};
  \node[proc, below=of graph] (scc) {Định kỳ: Tarjan SCC\\(phát hiện chu trình khả nghi)};

  \draw[flow] (e) -- (lpm);
  \draw[flow] (lpm) -- (cms);
  \draw[flow] (cms) -- (chk);
  \draw[flow] (chk) -- node[above, font=\scriptsize]{đúng} (al);
  \draw[flow] (chk) -- node[right, font=\scriptsize]{sai} (graph);
  \draw[flow] (al) |- (graph);
  \draw[flow] (graph) -- (scc);
\end{tikzpicture}
}
  \caption{Lưu đồ xử lý một sự kiện trong pipeline.}
  \label{fig:detect-flow-impl}
\end{figure}

Ở đây tôi đặt cả phân tích \emph{tần suất} (CMS) lẫn phân tích \emph{đồ thị}
(Tarjan SCC, BFS bán kính ảnh hưởng) \emph{ngay trong pipeline mỗi node}, chứ
không phải ở tầng tổng hợp. Tầng tổng hợp (mục~\ref{sec:impl-soc}) chỉ làm tương quan
và khử trùng lặp cảnh báo trên phạm vi toàn cụm.

\subsection{Phát cảnh báo}
\label{subsec:impl-pipeline-alert}

Cảnh báo được gửi qua một \emph{chuỗi xử lý} (\texttt{AlertHandlerChain}) cấu hình
được: handler ghi tệp \texttt{alerts.json} (NDJSON) và handler ghi \texttt{stderr}.
Mỗi cảnh báo mang đủ ngữ cảnh để SOC xử lý về sau: loại, nguồn, mức nghiêm trọng,
\texttt{node\_name}, \texttt{namespace}, nhãn và phần \texttt{details} dạng JSON (ví dụ
tần suất ước lượng, ngưỡng, danh sách đích trong bán kính ảnh hưởng). Mã
nguồn~\ref{lst:pipeline-emit} trích đoạn phát cảnh báo \texttt{port\_scan}.

\begin{lstlisting}[language=C++,caption={Phát cảnh báo \texttt{port\_scan} khi vượt ngưỡng (rút gọn từ \texttt{engine.cpp}).},label={lst:pipeline-emit}]
if (freq > threshold && alerted_sources.find(evt.src_ip) == alerted_sources.end()) {
    alerted_sources.insert(evt.src_ip);
    dsa::soc::Alert alert;
    alert.type      = "port_scan";
    alert.source    = ip_to_string(evt.src_ip);
    alert.severity  = severity_for(alert.type, cfg.severity_rules);
    alert.node_name = node_name;       // tu Downward API (NODE_NAME)
    alert.tags      = {"port_scan"};
    chain.send(alert);                 // -> alerts.json -> alert-bridge -> NATS
}
\end{lstlisting}

Container sidecar \texttt{alert-bridge} (Go) theo dõi đuôi tệp \texttt{alerts.json},
mỗi dòng mới được nó publish lên chủ đề NATS \texttt{zt.alerts.raw}. Tách phần đẩy NATS
ra một tiến trình Go riêng giúp bộ máy C++ không phải liên kết \texttt{libnats}, đồng
thời cho phép siết quyền sidecar tối đa (phân tích ở mục~\ref{sec:eval-trust}).

\section{Trung tâm vận hành an ninh (SOC)}
\label{sec:impl-soc}

SOC nhận cảnh báo thô từ các node, tương quan thành sự cố, làm giàu ngữ cảnh
Kubernetes, sinh chính sách cách ly và phục vụ người vận hành. Dòng chảy thông điệp:
pipeline mỗi node $\xrightarrow{\texttt{zt.alerts.raw}}$ \emph{aggregator}
$\xrightarrow{\texttt{zt.alerts.enriched}}$ \emph{incident-service}
$\rightarrow$ PostgreSQL $\rightarrow$ \emph{web-console}.

\subsection{Bộ tổng hợp (aggregator)}
\label{subsec:impl-aggregator}

Bộ tổng hợp (Hình~\ref{fig:comp-agg-bridge}, trái) đăng ký \texttt{zt.alerts.raw} theo
\emph{queue group} (để chạy nhiều bản sao mà không trùng lặp), đưa mỗi cảnh báo qua bộ
\emph{tương quan} (\texttt{Correlator}) rồi tái phát những cảnh báo đáng phát lên
\texttt{zt.alerts.enriched}. Khóa tương quan là \texttt{incident\_key = type|source};
hai cảnh báo cùng khóa trong cửa sổ thời gian được coi là cùng một sự cố. Logic cốt
lõi (Mã nguồn~\ref{lst:agg-correlate}): cảnh báo \emph{đầu tiên} của một khóa được phát
ngay; các cảnh báo lặp bị \emph{nén}; nhưng đúng thời điểm cùng một nguồn xuất hiện
trên node thứ hai (\emph{cross-node}), mức nghiêm trọng được \emph{nâng cấp} và cảnh
báo được tái phát kèm nhãn \texttt{cross\_node\_attacker}. Đây chính là góc nhìn
``toàn cụm'' của hệ thống~--~đạt được bằng tương quan cảnh báo, không phải bằng một đồ
thị toàn cục.

\begin{lstlisting}[language=Go,caption={Tương quan và nâng cấp cross-node trong aggregator (rút gọn từ \texttt{correlator.go}).},label={lst:agg-correlate}]
key := a.Type + "|" + a.Source
st, exists := c.incidents[key]
if !exists || now-st.LastSeenNS > c.Window.Nanoseconds() {
    st = &incidentState{Key: key, FirstSeenNS: now, LastSeenNS: now, Count: 1, ...}
    c.incidents[key] = st
    return a, true                       // lan dau -> phat ngay
}
prevCrossNode := len(st.Nodes) >= 2
st.LastSeenNS = now; st.Count++; st.Nodes[a.NodeName] = struct{}{}
if !prevCrossNode && len(st.Nodes) >= 2 {  // vua thanh cross-node
    a.Severity = upgradeSeverity(a.Severity)
    a.Tags = appendUnique(a.Tags, "cross_node_attacker")
    return a, true                       // nang cap -> tai phat
}
return nil, false                        // trung lap -> nen
\end{lstlisting}

\subsection{Dịch vụ sự cố (incident-service)}
\label{subsec:impl-incident}

Dịch vụ sự cố (Hình~\ref{fig:comp-inc-console}, trái) là thành phần phức tạp nhất của
SOC, gồm năm gói nội bộ: \texttt{natssub} (tiêu thụ cảnh báo), \texttt{store}
(PostgreSQL), \texttt{response} (sinh \texttt{NetworkPolicy} và thao tác pod),
\texttt{api} (REST) và \texttt{rbac}/\texttt{audit} (phân quyền và kiểm toán).

\paragraph{Tiêu thụ cảnh báo và làm giàu.} Bộ tiêu thụ (Mã nguồn~\ref{lst:inc-sub})
\emph{upsert} mỗi cảnh báo enriched thành sự cố theo \texttt{incident\_key}: nếu chưa
có thì tạo mới, nếu có thì cập nhật \texttt{last\_seen}/số lần/nodes/mức nghiêm trọng.
Với sự cố mới, tôi cho dịch vụ thực hiện hai việc: (i)~\emph{làm giàu ngữ cảnh
Kubernetes}~--~tra cứu pod nào đang giữ địa chỉ nguồn (\texttt{LookupPodByIP} duyệt
\texttt{status.PodIP}) rồi lưu kèm tên pod/namespace, để người vận hành thấy ngay danh
tính kẻ tấn công thay vì chỉ một địa chỉ IP; (ii)~\emph{tự sinh hành động} khi mức
nghiêm trọng thuộc \{high, critical\}~--~tạo bản nháp \texttt{NetworkPolicy} ở trạng
thái \texttt{pending} chờ duyệt.

\begin{lstlisting}[language=Go,caption={Tiêu thụ cảnh báo: upsert sự cố, làm giàu pod, tự sinh hành động (rút gọn từ \texttt{subscriber.go}).},label={lst:inc-sub}]
inc, created, _ := st.UpsertIncident(ctx, &inAlert)
st.AppendAlert(ctx, inc.ID, msg.Data)
if created && resolver != nil {                       // lam giau ngu canh K8s
    if name, found, _ := resolver.LookupPodByIP(ctx, ns, inc.SourceIP); found {
        st.SetIncidentPodContext(ctx, inc.ID, name, ns)
    }
}
if _, ok := autoSet[inc.Severity]; ok {               // high|critical -> sinh nhap chinh sach
    yaml := response.GenerateQuarantine(response.QuarantineRequest{
        IncidentID: inc.ID, SourceIP: inc.SourceIP, Severity: inc.Severity})
    st.CreateAction(ctx, &store.Action{IncidentID: inc.ID, Type: "quarantine_pod", YAMLDraft: yaml})
}
\end{lstlisting}

\paragraph{Lưu trữ.} Trạng thái được giữ trong PostgreSQL qua bốn bảng~--~\texttt{incidents},
\texttt{alerts}, \texttt{actions}, \texttt{audit\_log}~--~mà lược đồ chi tiết đã trình
bày ở ERD Chương~\ref{ch:design}. Tôi gộp các lần xuất hiện lặp của cùng một sự cố bằng
phép \emph{upsert} \texttt{INSERT ... ON CONFLICT (incident\_key) DO UPDATE}, dựa trên
ràng buộc duy nhất trên \texttt{incident\_key} (Mã nguồn~\ref{lst:sql-upsert}); migration
được nhúng vào nhị phân và áp tự động khi khởi động.

\begin{lstlisting}[language=SQL,caption={Lược đồ bảng \texttt{incidents} và khóa duy nhất dùng cho upsert (\texttt{001\_init.sql}).},label={lst:sql-upsert}]
CREATE TABLE IF NOT EXISTS incidents (
    id               BIGSERIAL PRIMARY KEY,
    incident_key     TEXT NOT NULL,                 -- type|source: khoa khu trung lap
    type             TEXT NOT NULL,
    source_ip        TEXT NOT NULL,
    severity         TEXT NOT NULL DEFAULT 'medium',
    status           TEXT NOT NULL DEFAULT 'new',   -- new|ack|resolved|suppressed
    nodes_observed   TEXT[] NOT NULL DEFAULT '{}',
    occurrence_count INT  NOT NULL DEFAULT 1,
    cross_node       BOOLEAN NOT NULL DEFAULT FALSE,
    created_at       TIMESTAMPTZ NOT NULL DEFAULT NOW(),
    UNIQUE (incident_key)
);
\end{lstlisting}

\paragraph{Sinh và áp chính sách cách ly.} Bộ sinh phản hồi (\texttt{response}) tạo
\texttt{NetworkPolicy} dạng \emph{chặn mọi lưu lượng đi ra} (egress-deny) gắn lên đúng
pod tấn công qua một nhãn theo sự cố. Khi hiện thực bước này tôi gặp một vấn đề kỹ
thuật quan trọng: vì ngữ nghĩa \texttt{NetworkPolicy} là \emph{cộng dồn}
(một kết nối được phép nếu \emph{bất kỳ} policy nào cho phép), nên kiểu ``cho phép mọi
nguồn trừ kẻ tấn công'' sẽ bị các luật cho-phép rộng khác vô hiệu hóa; ngược lại, chặn
\emph{egress của chính pod tấn công} thì không policy ingress nào ghi đè được, và dừng
được cả quét cổng lẫn dịch chuyển ngang tại nguồn. Cơ chế áp dụng theo
\emph{có người duyệt}: khi admin phê duyệt, hệ thống mới gắn nhãn cách ly lên pod
(\texttt{QuarantinePodByIP}) rồi áp policy qua API Kubernetes. Toàn bộ luồng này được
tóm tắt ở Hình~\ref{fig:flow-netpol}, còn bản nháp sinh ra có dạng Mã
nguồn~\ref{lst:netpol}.

\begin{figure}[H]
  \centering
  \resizebox{0.92\linewidth}{!}{% Luu do sinh & ap NetworkPolicy (generate-for-review). Chuong 4.
\begin{tikzpicture}[font=\footnotesize, >={Stealth[length=2.2mm]}, node distance=0.62cm,
    io/.style  ={draw, rounded corners=7pt, fill=blue!10, align=center,
                 minimum width=5.2cm, minimum height=0.66cm},
    proc/.style={draw, fill=gray!8, align=center, minimum width=5.2cm, minimum height=0.66cm},
    gw/.style  ={draw, diamond, aspect=2.2, fill=yellow!22, align=center,
                 inner sep=1pt, minimum width=3.2cm},
    term/.style={draw, rounded corners=7pt, fill=green!14, align=center,
                 minimum width=2.8cm, minimum height=0.6cm},
    flow/.style={->, thick},
  ]
  \node[io]   (alert)  {Cảnh báo \texttt{port\_scan} (qua NATS)};
  \node[proc, below=of alert]  (inc)    {Tạo sự cố (\texttt{incident-service})};
  \node[proc, below=of inc]    (enrich) {Làm giàu: \texttt{src\_ip} $\rightarrow$ pod (tên, namespace)};
  \node[proc, below=of enrich] (gen)    {Sinh bản nháp \texttt{NetworkPolicy} egress-deny\\gắn nhãn cách ly lên đúng pod};
  \node[gw,   below=0.7cm of gen] (appr) {Người vận hành\\duyệt?};
  \node[proc, below=0.78cm of appr] (apply)  {Áp policy qua K8s API (Calico thực thi)};
  \node[proc, below=of apply]  (verify) {Xác minh: pod tấn công mất kết nối đi ra};
  \node[term, below=of verify] (done)   {Ghi nhật ký kiểm toán};
  \node[proc, right=1.5cm of appr] (reject) {Ghi kiểm toán \& tiếp tục theo dõi};

  \draw[flow] (alert)--(inc);   \draw[flow] (inc)--(enrich);
  \draw[flow] (enrich)--(gen);  \draw[flow] (gen)--(appr);
  \draw[flow] (appr)-- node[right]{đồng ý} (apply);
  \draw[flow] (apply)--(verify); \draw[flow] (verify)--(done);
  \draw[flow] (appr)-- node[above]{từ chối} (reject);

  % Chu thich ky thuat ve egress-deny
  \node[draw, dashed, fill=red!4, align=left, font=\scriptsize, text width=4.0cm,
        left=1.4cm of gen] (note)
    {\textbf{Vì sao egress-deny?} Ngữ nghĩa \emph{cộng dồn} của NetworkPolicy khiến luật
     ``cho phép mọi nguồn trừ một'' không chặn được khi còn luật cho-phép rộng khác; phải
     chặn \emph{lối ra} của chính pod tấn công.};
  \draw[->, red!60, dashed] (note) -- (gen);
\end{tikzpicture}
}
  \caption{Lưu đồ sinh và áp \texttt{NetworkPolicy} cách ly.}
  \label{fig:flow-netpol}
\end{figure}

\begin{lstlisting}[language={},caption={Bản nháp \texttt{NetworkPolicy} egress-deny sinh tự động cho pod tấn công.},label={lst:netpol}]
apiVersion: networking.k8s.io/v1
kind: NetworkPolicy
metadata:
  name: zt-quarantine-incident-<id>
  namespace: zt-targets
spec:
  podSelector:
    matchLabels:
      zt-soc-quarantine-<id>: "true"   # pod tan cong, dan nhan boi enforcer
  policyTypes: [Egress]
  egress: []                           # chan TOAN BO luu luong di ra cua pod do
\end{lstlisting}

\paragraph{API, phân quyền và kiểm toán.} Dịch vụ phơi bày REST API (router
\texttt{chi}) với phân quyền theo vai trò ba mức~--~\texttt{viewer} (chỉ đọc),
\texttt{editor} (ghi nhận/đóng sự cố), \texttt{admin} (phê duyệt/bác bỏ phản hồi). Phân
quyền hiện thực bằng middleware xác thực Bearer token và bộ chặn \texttt{RequireRole}
theo bậc thang \texttt{admin > editor > viewer} (Mã nguồn~\ref{lst:rbac}). Mọi thao
tác làm thay đổi trạng thái~--~tạo sự cố, đổi trạng thái, sinh/duyệt/bác bỏ hành
động~--~đều được ghi vào bảng \texttt{audit\_log} kèm chủ thể, vai trò và phần
\texttt{diff}, phục vụ truy vết và xuất báo cáo tuân thủ.

\begin{lstlisting}[language=Go,caption={Định tuyến có phân quyền: phê duyệt hành động yêu cầu vai trò \texttt{admin} (rút gọn từ \texttt{api.go}, \texttt{rbac.go}).},label={lst:rbac}]
r.Use(rbac.Middleware(tokens))                                  // xac thuc Bearer token
r.Route("/api/v1/actions", func(r chi.Router) {
    r.Get("/", s.listActions)                                   // viewer doc duoc
    r.With(rbac.RequireRole(rbac.RoleAdmin)).Post("/{id}/approve", s.approveAction)
    r.With(rbac.RequireRole(rbac.RoleAdmin)).Post("/{id}/reject",  s.rejectAction)
})
// RequireRole: admin(3) > editor(2) > viewer(1)
func RequireRole(min Role) func(http.Handler) http.Handler { /* tu choi neu rank < min */ }
\end{lstlisting}

\subsection{Console web}
\label{subsec:impl-console}

Console web (Hình~\ref{fig:comp-inc-console}, phải) là một giao diện mỏng Flask~+~HTMX
đặt trước REST API của dịch vụ sự cố. Nó liệt kê sự cố và hành động đang chờ, cho phép
admin phê duyệt/bác bỏ \texttt{NetworkPolicy} ngay tại dòng (HTMX thay từng hàng mà
không tải lại trang), và xem báo cáo tuần. Về bảo mật, console xác thực người dùng bằng
Basic Auth, còn khi gọi xuống API thì dùng \emph{token admin phía máy chủ}
(\texttt{RBAC\_ADMIN\_TOKEN})~--~trình duyệt không bao giờ giữ token API. Đây là một
proxy có chủ đích: tách phần xác thực người dùng khỏi phần phân quyền API.

\section{Ứng dụng và giao diện vận hành (chạy thật)}
\label{sec:impl-app}

Mục này trình bày ứng dụng \emph{chạy thật} trên cụm, phân biệt với phần thiết kế giao
diện (mockup) ở Chương~\ref{ch:design}. Console phục vụ trọn vòng xử lý sự cố: người
vận hành thấy danh sách sự cố có thể lọc/sắp xếp, mở chi tiết một sự cố để xem siêu dữ
liệu đã làm giàu (gồm pod nguồn), các cảnh báo liên quan và bản nháp chính sách cách
ly, rồi ghi nhận hoặc phê duyệt tùy vai trò. Song song, dashboard Grafana cung cấp góc
nhìn kỹ thuật theo thời gian thực về số lượng cảnh báo, dòng thời gian sự cố và các
nguồn hoạt động mạnh nhất.

Console được phục vụ tại \texttt{console.zt.local} qua ingress của cụm. Màn hình
\emph{Pending actions} (Hình~\ref{fig:console-actions}) là hàng đợi phê duyệt: mỗi dòng
là một bản nháp hành động (chủ yếu \texttt{quarantine\_pod}) do dịch vụ sự cố tự sinh từ
cảnh báo mức cao, kèm hai nút \emph{Approve}/\emph{Reject}~--~hiện thực hóa điểm dừng
\emph{có người trong vòng lặp}. Tại thời điểm chụp, console liệt kê $38$ hành động đang
chờ, đúng bằng số sự cố mức \texttt{high}/\texttt{critical} đang lưu trong cơ sở dữ liệu.

\begin{figure}[H]
  \centering
  \includegraphics[width=\linewidth]{console-actions}
  \caption{Hàng đợi phê duyệt hành động trên console.}
  \label{fig:console-actions}
\end{figure}

Mở chi tiết một sự cố (Hình~\ref{fig:incident-detail}) cho thấy toàn bộ ngữ cảnh đã làm
giàu: địa chỉ nguồn, mức nghiêm trọng, các node quan sát được, cờ \emph{cross-node}, số
lần tái phát, mốc thời gian đầu/cuối, danh sách hành động, biểu mẫu thao tác tầng~2 và
\emph{nhật ký kiểm toán}. Ví dụ sự cố \texttt{blast\_radius} mức \texttt{critical} từ
nguồn \texttt{100.119.246.28} được quan sát đồng thời trên hai node
(\texttt{userver-master}, \texttt{userver-worker-home}) nên được nâng mức và đánh dấu
cross-node.

\begin{figure}[H]
  \centering
  \includegraphics[width=\linewidth]{incident-detail}
  \caption{Chi tiết một sự cố với ngữ cảnh đã làm giàu và nhật ký kiểm toán.}
  \label{fig:incident-detail}
\end{figure}

Sau khi quản trị viên phê duyệt, hành động chuyển trạng thái \texttt{executed} và đối
tượng \texttt{NetworkPolicy} tương ứng được áp vào cụm~--~Hình~\ref{fig:incident-enforced}
ghi lại dòng \texttt{networkpolicy.networking.k8s.io/zt-quarantine-incident-4} đã được áp
trong không gian tên \texttt{zt-targets}, minh chứng vòng phản hồi đã khép kín tới tầng
thực thi.

\begin{figure}[H]
  \centering
  \includegraphics[width=\linewidth]{incident-enforced}
  \caption{Hành động cách ly đã được áp (\texttt{executed}) sau phê duyệt.}
  \label{fig:incident-enforced}
\end{figure}

Trang \emph{Reports} tổng hợp báo cáo tuần đọc trực tiếp từ PostgreSQL
(Hình~\ref{fig:weekly-report}): tiêu đề kỳ, tổng số sự cố, phân bố theo mức nghiêm trọng
và bảng nguồn tấn công nhiều nhất, kèm nút \emph{Generate \& upload to MinIO} để kết
xuất và lưu trữ báo cáo. Đây cũng là nơi thể hiện rõ rằng số liệu hiển thị trên
console/Reports là \emph{số liệu thật} tích lũy trong cơ sở dữ liệu, không phải dữ liệu
dựng sẵn.

\begin{figure}[H]
  \centering
  \includegraphics[width=\linewidth]{weekly-report}
  \caption{Trang báo cáo tuần kết xuất từ cơ sở dữ liệu.}
  \label{fig:weekly-report}
\end{figure}

\paragraph{Về đường dữ liệu của dashboard Grafana.} Cần phân biệt hai nguồn dữ liệu
khác nhau. Console/Reports đọc \emph{trực tiếp} PostgreSQL~--~\emph{hệ thống lưu trữ bền
vững} tích lũy mọi sự cố từ trước tới nay (Hình~\ref{fig:weekly-report}). Trong khi đó
các dashboard Grafana là \emph{góc nhìn quan trắc theo cửa sổ thời gian}: bảng
\emph{Alerts Overview} truy vấn nhật ký trong Loki theo cửa sổ mặc định $30$~phút, còn
bảng \emph{Incidents \& Actions} truy vấn các bộ đếm Prometheus của dịch vụ SOC. Do đó
dashboard chỉ phản ánh hoạt động \emph{gần đây} và đòi hỏi (i)~có lưu lượng tấn công
trong cửa sổ đang xem, và (ii)~các bộ đếm Prometheus phải được thu thập. Hệ quả của sự
khác biệt nguồn dữ liệu này~--~và cấu hình thu thập tương ứng~--~được bàn kỹ hơn ở
mục giới hạn của Chương~\ref{ch:eval}.

Sau khi khai báo \texttt{ServiceMonitor} cho các dịch vụ SOC, Prometheus thu thập đúng
hai mục tiêu \texttt{zt-aggregator:9100} và \texttt{zt-incident-service:9101}; hai
dashboard \emph{ZT-SOC Overview} và \emph{ZT-SOC Incidents}
(Hình~\ref{fig:grafana-overview} và~\ref{fig:grafana-incidents}) hiển thị các chỉ số như
số sự cố đang hoạt động, sự cố cross-node, tốc độ nhận cảnh báo và tốc độ tạo sự cố. Do
các panel dạng \emph{rate} phản ánh \emph{biến thiên} của bộ đếm, đường biểu diễn chỉ
nhô lên khi có lưu lượng tấn công \emph{ngay trong cửa sổ thời gian đang xem}; vì vậy ảnh
chụp được lấy bằng cách thu hẹp cửa sổ (vài phút), bật tự làm mới và chạy kịch bản tấn
công (Phụ lục~\ref{app:demo-attack}) ngay trước khi quan sát.

\begin{figure}[H]
  \centering
  \IfFileExists{assets/images/grafana-overview.png}{%
    \includegraphics[width=\linewidth]{grafana-overview}}{%
    \fbox{\parbox[c][2.6cm][c]{0.8\linewidth}{\centering\itshape\small
    [Ảnh dashboard \emph{ZT-SOC Overview} --- lưu vào
    \texttt{stage2/assets/images/grafana-overview.png}]}}}
  \caption{Dashboard \emph{ZT-SOC Overview} trên Grafana.}
  \label{fig:grafana-overview}
\end{figure}

\begin{figure}[H]
  \centering
  \IfFileExists{assets/images/grafana-incidents.png}{%
    \includegraphics[width=\linewidth]{grafana-incidents}}{%
    \fbox{\parbox[c][2.6cm][c]{0.8\linewidth}{\centering\itshape\small
    [Ảnh dashboard \emph{ZT-SOC Incidents} --- lưu vào
    \texttt{stage2/assets/images/grafana-incidents.png}]}}}
  \caption{Dashboard \emph{ZT-SOC Incidents} trên Grafana.}
  \label{fig:grafana-incidents}
\end{figure}

\section{Chạy thử nghiệm hệ thống end-to-end}
\label{sec:impl-e2e}

Để kiểm chứng toàn bộ vòng phát hiện~$\rightarrow$~phản hồi hoạt động liền mạch, hệ
thống được chạy thử theo một kịch bản end-to-end trên cụm thật.
Hình~\ref{fig:flow-e2e} tóm tắt kịch bản qua các thành phần đã triển khai, từ pod tấn
công tới bước xác minh cách ly.

\begin{figure}[H]
  \centering
  \resizebox{0.62\linewidth}{!}{% Luu do kich ban end-to-end: attacker -> phat hien -> su co -> cach ly -> xac minh. Chuong 4.
\begin{tikzpicture}[font=\footnotesize, >={Stealth[length=2.4mm]}, node distance=0.6cm,
    comp/.style={draw, rounded corners=2pt, align=center, minimum width=6.6cm,
                 minimum height=0.78cm},
    atk/.style ={comp, fill=red!14},
    kern/.style={comp, fill=green!12},
    proc/.style={comp, fill=blue!10},
    soc/.style ={comp, fill=orange!14},
    hum/.style ={comp, fill=violet!12},
    flow/.style={->, very thick},
    lbl/.style ={font=\scriptsize\itshape, fill=white, inner sep=1pt},
  ]
  \node[atk]  (atk)    {Pod tấn công --- quét cổng \texttt{nmap -sT}};
  \node[kern, below=of atk]    (col) {\texttt{collector-ebpf} (DaemonSet): kprobe \texttt{tcp\_v4\_connect},\\lọc LPM (chỉ mạng pod), đẩy qua ring buffer};
  \node[proc, below=of col]    (pipe){\texttt{pipeline}: CMS \texttt{record(src\_ip)}; vượt ngưỡng $40$/$60$s};
  \node[soc,  below=of pipe]   (inc) {\texttt{incident-service}: tạo sự cố $+$ làm giàu \texttt{src\_ip}$\rightarrow$pod;\\sinh bản nháp \texttt{NetworkPolicy} egress-deny};
  \node[hum,  below=of inc]    (con) {\texttt{web-console}: người vận hành xem xét \& \emph{phê duyệt}};
  \node[proc, below=of con]    (k8s) {K8s API $+$ Calico: áp \texttt{NetworkPolicy} cách ly};
  \node[kern, below=of k8s]    (ver) {Xác minh: pod tấn công mất egress; pod hợp lệ không ảnh hưởng};

  \draw[flow] (atk)  -- node[lbl,right]{\texttt{connect()}} (col);
  \draw[flow] (col)  -- node[lbl,right]{dòng JSON} (pipe);
  \draw[flow] (pipe) -- node[lbl,right]{cảnh báo \texttt{port\_scan} (NATS)} (inc);
  \draw[flow] (inc)  -- node[lbl,right]{bản nháp egress-deny} (con);
  \draw[flow] (con)  -- node[lbl,right]{\texttt{apply} (sau phê duyệt)} (k8s);
  \draw[flow] (k8s)  -- node[lbl,right]{thực thi ở tầng mạng} (ver);
\end{tikzpicture}
}
  \caption{Lưu đồ kịch bản thử nghiệm end-to-end.}
  \label{fig:flow-e2e}
\end{figure}

Một pod đóng vai kẻ tấn công được đặt trên một node worker và quét cổng bằng
\texttt{nmap} ở chế độ connect (\texttt{-sT}) nhắm vào các pod dịch vụ. Bộ thu thập
eBPF bắt các lời gọi \texttt{connect()} và đẩy sự kiện lên pipeline; khi số kết nối từ
pod tấn công vượt ngưỡng trong cửa sổ thời gian, pipeline phát cảnh báo
\texttt{port\_scan} (và \texttt{blast\_radius} khi pod chạm tới nhiều đích) lên NATS.
Bộ tổng hợp khử trùng lặp và nâng cấp mức nghiêm trọng nếu nguồn xuất hiện trên nhiều
node; dịch vụ sự cố tạo sự cố, làm giàu địa chỉ nguồn về đúng pod tấn công, rồi sinh
bản nháp \texttt{NetworkPolicy} chặn mọi lưu lượng đi ra của pod đó. Người vận hành xem
xét và phê duyệt qua console (điểm dừng có người trong vòng lặp); chỉ khi đó chính sách
mới được áp vào cụm qua API Kubernetes. Bước cuối xác nhận hiệu lực: sau khi áp chính
sách, pod tấn công không còn thiết lập được kết nối đi ra, trong khi lưu lượng hợp lệ
của các pod khác không bị ảnh hưởng.

Mã nguồn~\ref{lst:e2e-log} cho thấy chính xác hệ thống chạy ra như thế nào trong một
lần thử: nhật ký của pipeline ghi cảnh báo \texttt{port\_scan}/\texttt{blast\_radius},
dịch vụ sự cố tạo sự cố và tự sinh hành động cách ly, và sau khi tôi phê duyệt trên
console thì \texttt{NetworkPolicy} được áp khiến pod tấn công mất kết nối đi ra (mã trả
về của lệnh kết nối đổi từ $0$ sang khác $0$).

\begin{lstlisting}[language={},caption={Nhật ký một lần chạy thử end-to-end (rút gọn từ sự cố \#4 thật trên cụm).},label={lst:e2e-log}]
# 1) Pipeline phat hien quet cong + blast radius (node master)
$ kubectl -n zt-mapper logs zt-pipeline-ebpf-qztgt -c pipeline | grep ALERT
[ALERT] Port scan detected: 100.119.246.28 (freq > tau=40)
[ALERT] Blast radius from 100.119.246.28: multi-node within 3 hops

# 2) Dich vu su co tao su co cross-node + tu sinh ban nhap NetworkPolicy
$ kubectl -n soc logs deploy/zt-incident-service | grep -E 'created|action'
[sub] incident.created id=4 type=blast_radius source=100.119.246.28 severity=critical
[sub] action.auto_generated incident_id=4 type=quarantine_pod

# 3) Sau khi admin phe duyet tren console -> ap chinh sach cach ly
[apply] networkpolicy.networking.k8s.io/zt-quarantine-incident-4 created -n zt-targets

# 4) Xac minh: pod tan cong mat ket noi di ra (truoc cach ly exit=0)
$ kubectl -n zt-targets exec attacker -- nc -w2 10.244.204.161 80; echo "exit=$?"
exit=1
\end{lstlisting}

Kịch bản này chứng minh tính \emph{đúng đắn về mặt chức năng} của toàn hệ thống; các số
liệu định lượng (tỷ lệ phát hiện, FPR, chi phí, độ trễ) được phân tích ở
Chương~\ref{ch:eval}. Mã nguồn đầy đủ của các thành phần cốt lõi được đưa vào phần Phụ
lục: chương trình eBPF (Phụ lục~\ref{app:ebpf}), console web (Phụ
lục~\ref{app:web-console}) và lược đồ dữ liệu PostgreSQL (Phụ lục~\ref{app:sql}).

\section{Kết luận chương}

Chương này đã trình bày trọn vẹn quá trình hiện thực và vận hành hệ thống trên một cụm
Kubernetes thật, bao trùm cả ba lớp: hạ tầng nền tảng theo mô hình GitOps (giới thiệu
vắn tắt), ba thành phần lõi được đào sâu kèm trích đoạn mã nguồn/cấu hình thật~--~bộ
thu thập eBPF, pipeline phát hiện và các dịch vụ SOC~--~cùng ứng dụng giao diện cho
người vận hành. Riêng \emph{mô hình tin cậy} của bộ thu thập đặc quyền được phân tích
cùng các giới hạn an ninh ở Chương~\ref{ch:eval}.
Việc chạy thử end-to-end cho thấy hệ thống không dừng ở mức nguyên mẫu rời rạc mà thực
sự khép kín được vòng phát hiện~$\rightarrow$~phản hồi với con người ở đúng điểm ra
quyết định. Đây là tiền đề trực tiếp cho phần đánh giá định lượng ở chương tiếp theo,
nơi từng mắt xích của vòng này được đo đạc và phân tích trên dữ liệu thật.
